\chapter{Conclusions and Future Work}
\label{ch:conclusions-future-work}

The work presented in this thesis achieved its scoped objective: it produced
and evaluated a node-local runtime security monitoring prototype for the
selected Rocky Linux environment. The result is not a general security
platform, but an end-to-end implementation that connects bounded kernel
observations to typed evidence and deterministic local detection. The
conclusions that follow are therefore grounded in the implemented event domain
and in the controlled evaluation reported in
Chapter~\ref{ch:experimental-evaluation}.

\section{Summary of the Work}
\label{sec:conclusions-summary-work}

The central problem was to collect heterogeneous process and security-related
kernel observations without making the collection layer responsible for their
complete interpretation. The resulting architecture assigns constrained data
extraction and optional early admission to \gls{ebpf}, while the Go runtime
owns record normalisation, monitoring scope, detection and presentation. A
shared event contract connects these responsibilities so that observations
from different physical hooks can enter one logical event model.

In kernel space, the prototype combines complementary observation points to
retain different stages of an operation, including requests, security
mediation, committed transitions and returned outcomes. Direct and paired
producers use a common bounded construction framework, selective probe
planning instead activates only the producers and support programs required by the
chosen logical events. Exact \gls{uid} admission can reject observations before
full enrichment and transport when the monitored subject has a suitable fixed
identity. Compatibility work based on CO-RE, explicit layout branches and
verifier-visible bounds allows these mechanisms to operate on the verified
Rocky Linux kernel.

In user space, schema-directed decoding converts perf-buffer records into the
typed event representation shared by later components. Policies define the
relevant event domain and suppression decisions, detectors instead evaluate
the admitted evidence. Point detectors act on one event, and collective
detectors retain bounded local state for short ordered sequences. Alerts
preserve the observations responsible for a match and may propagate
\gls{attack} tactics and techniques as classification metadata. Together,
these components form an engineering integration contribution: established
kernel instrumentation and rule-based detection mechanisms were combined into
a compact pipeline adapted to a constrained enterprise-kernel target.

\section{Main Results}
\label{sec:conclusions-main-results}

The functional campaign established the intended node-local observation
boundary. In the single-node Minikube testbed, the monitor, vulnerable target
and trigger ran in three separate Pods. All three exploit repetitions recovered
the synthetic flag and produced one complete alert containing the configured
four-event sequence. The three repetitions of the selected benign interaction
produced the expected application response without completing that alert. The
result demonstrates repeatable same-node, cross-Pod visibility and detection
for this controlled scenario. In particular, the experiment exercised the
same-process path of the selected collective detector rather than its complete
set of possible process relationships.

The performance comparison also supports the value of narrow early admission
under the measured conditions. Across three runs, the exact UID-filtered
profile consumed a mean of 0.46\% of one \gls{cpu} core in the monitor-container
cgroup, and every filtered run remained below the project's 5\% gate. The
corresponding unfiltered mean was 5.73\%; using the retained data, the filtered
mean was approximately 92.0\% lower while flag recovery and the complete alert
were preserved in every monitored run. This measurement concerns CPU charged
to the monitor cgroup as a percentage of one core. It excludes kernel work
charged elsewhere and must not be interpreted as total eBPF, node or system
overhead. It nevertheless provides scenario-specific evidence that discarding
irrelevant observations before transport can substantially reduce the demand
placed on the user-space monitoring path.

\section{Limitations}
\label{sec:conclusions-limitations}

The prototype's technical assurance remains narrower than its implemented
event surface. Kernel and Go representations of the event contract are
maintained separately, so consistency depends on review and cross-registry
tests rather than generation from one source. The current registry gate also
exposes an unresolved classification inconsistency: an internal socket support
probe is treated as a public producer despite having no decoder schema.
Object-backed tests do not yet exercise every loading, attachment,
partial-initialisation and cleanup path on the target kernel. External policy
and detector documents are validated semantically, but unknown YAML fields are
not rejected strictly. Collective state is bounded, whereas some maintenance
and alert-deduplication behaviour still depends on subsequent traffic and does
not have equally strong capacity guarantees.

The analytical and operational scope is intentionally local. Detection is
deterministic and configuration driven; the system neither discovers unknown
behaviour through a learned baseline nor reconstructs persistent process or
cluster graphs. The prototype observes and reports activity but does not enforce 
decisions.

The evaluation further limits the breadth of the conclusions. It covered one
enterprise kernel, one two-core Minikube node, one vulnerable workload and one
scenario-specific collective detector, with three repetitions per condition.
Consequently, the campaign supports a transparent engineering demonstration, not production readiness, broad
portability, loss-free delivery or statistical claims about detection quality.

\section{Future Work}
\label{sec:conclusions-future-work}

The future development of the prototype should concentrate on extending its
  operational capacity while preserving the separation between bounded kernel
  collection and configurable user-space analysis. The principal priorities are
  scalability, performance, broader early filtering, additional host-level
  observation points and richer evidence from the hooks already implemented.

  Scalability should be considered at both agent and deployment level. A single
  instance must remain responsive as the number of enabled events, policies,
  detectors and monitored processes increases. This requires stress tests with
  progressively higher event rates, longer execution periods and larger detector
  catalogues. In a multi-node environment, each agent should continue to process
  local observations independently, while a versioned export interface could
  deliver events and alerts to an external platform for storage and
  cluster-level analysis. Such an architecture would preserve bounded local
  state and avoid introducing distributed correlation into the kernel-facing
  agent.

  Performance work should extend the measurements presented in
  Chapter~\ref{ch:experimental-evaluation}. Future experiments should account
  separately for kernel execution, event transport, decoding, policy evaluation,
  detector processing and output rendering. Per-event and per-probe counters
  would help identify the observation points responsible for the greatest cost.
  Known-volume workloads could also quantify transport loss and establish the
  event rates sustained by different perf-buffer capacities. Alternative
  transport mechanisms, allocation reduction and detector-state optimisation
  should be considered after profiling identifies the dominant bottlenecks.

  The current exact UID admission demonstrates that early filtering can reduce
  the amount of data transported to user space, although its applicability
  depends on prior knowledge of the monitored identity. A broader low-level
  filtering model could support additional bounded selectors, including host
  process identifiers, group identifiers, command names, cgroup identifiers,
  ecc.. Where practical, the control plane could translate simple policy constraints into kernel map
  entries and update them without rebuilding the eBPF object. Complex conditions
  and detector logic would remain in user space, where they can evolve without
  increasing verifier complexity or embedding high-level security semantics in
  the collectors.

  The host event domain can be expanded with new hooks selected from concrete
  security use cases. Relevant areas include additional file modification outcomes, 
  kernel key management, process isolation changes and other interfaces that can alter host behaviour. Each new
  logical event should identify the precise kernel stage being observed and
  distinguish requests, mediation points, committed transitions and returned
  outcomes. Hook availability and payload extraction must also be verified
  against the supported enterprise kernels before an event is included in the
  public catalogue.

  Existing hooks can provide greater analytical value through carefully selected
  additional fields. Useful extensions include return evidence for operations
  currently observed only at entry, richer capability-check context, stable parent identity, and more
  complete requested-versus-resolved resource details. These fields would allow
  security analysts to distinguish similar operations more precisely and write
  detectors with fewer ambiguous matches. Enrichment should remain bounded and
  should be introduced only when the field contributes directly to a detection
  or investigation requirement.

  Every extension affects the shared event contract and decoder schema. 
  New filters, hooks and fields should therefore be
  accompanied by registry checks, decoder fixtures, target-kernel attachment
  tests and workload-based measurements.